% 出席確認クイズ 第06回　組織・人材とAI協働
%   attendance/attendance06.html から生成（解説は本ガイド用に加筆）。

\quizq{1}{「人とAIの協働」の説明として適切なものはどれでしょうか?}%
  {\correct{AIの提案を人が確認・修正しながら成果を出す働き方}}%
  {AIがすべてを決め、人は関与しない}%
  {AI同士が競争する}%
  {人がAIをまったく使わない}%
  {正解は (a)。人とAIの協働では、AI の提案を人が確認・修正し、それぞれの得意分野を活かして成果を出します。(b) は全自動化であり、協働ではありません。AIを下書き担当・壁打ち相手・点検役のどれとして使うかを意識すると、協働の設計がしやすくなります。}

\quizq{2}{AI導入の際に組織で起こりやすい抵抗の要因として適切なものはどれでしょうか?}%
  {電力不足}%
  {\correct{仕事がなくなる不安や使い方への不慣れ}}%
  {為替の変動}%
  {天候の変化}%
  {正解は (b)。AI導入への抵抗は、雇用への不安と新しい道具への不慣れが主な要因です。(a)(c)(d) は外部環境の要因で、組織内部の抵抗とは関係がありません。抵抗は非合理ではなく、合理的な不安の表れとして理解する必要があります。}

\quizq{3}{変革管理(チェンジマネジメント)で重要なこととして適切なものはどれでしょうか?}%
  {罰則で強制する}%
  {一度にすべてを切り替える}%
  {\correct{目的の共有と段階的な導入、教育や支援}}%
  {導入を従業員に秘密にする}%
  {正解は (c)。変革を定着させるには、目的の共有、段階的な導入、教育・支援による不安の解消が重要です。(a)(b)(d) はいずれも抵抗を強めます。コッターの変革8段階（危機感の醸成、ビジョンの共有、短期的成果など）も参考になります。}

\quizq{4}{生成AIの普及によって重要性が高まると考えられるスキルはどれでしょうか?}%
  {電話応対のみ}%
  {\correct{適切な問いを立て、AIの出力を評価する力}}%
  {暗算の速さ}%
  {手書きの速さ}%
  {正解は (b)。AIが文章や案を出せる時代には、適切な問いを設計し、出力の質を評価・判断する力がより重要になります。(a)(c)(d) は機械に置き換えられやすい技能です。評価には対象領域の専門知識が必要で、専門性の価値がなくなるわけではありません。}

\quizq{5}{AI導入において経営層に求められる役割として適切なものはどれでしょうか?}%
  {現場に任せて一切関与しない}%
  {技術の詳細をすべて自分で実装する}%
  {\correct{目的と方針を示し、責任を持って推進する}}%
  {AIの利用をすべて禁止する}%
  {正解は (c)。経営層には、AI活用の目的と方針を示し、体制や投資を決め、責任を持って推進する役割があります。(b) の技術実装は専門家の仕事で、(a) の現場への丸投げは第7回で扱う失敗要因の1つです。(d) は競争力の低下とシャドーAIを招きます。}
